%-----------------------------------------------------------------------------
% Shared notation and reference helpers for PDF, lwarp, and TeX4ht.
% Loaded by each backend after its package and layout setup, before ntheorem.
% Keep rendering-dependent definitions in the corresponding backend prefix.
% Add new backend-independent notation here so all outputs share it.
%-----------------------------------------------------------------------------

\providecommand{\remove}[1]{}
\providecommand{\NotAMSMode}[1]{#1}
\providecommand{\Solution}[1]{}
\providecommand{\Version}[1]{\hfill{Version: #1}}
\providecommand{\Distor}{\mathrm{dist}}
\providecommand{\Hardness}[1]{(h:#1)}
\providecommand{\QuoteOptExt}[3]{\QuoteExt{#1}{#2}{#3}}
\providecommand{\Hypercube}{\EuScript{H}}
\providecommand{\CoverTime}{\EuScript{C}}
\providecommand{\Nbr}{\Gamma}
\providecommand{\CommuteTime}[2]{\mathbf{CT}_{#1#2}}
\providecommand{\Resistance}{\mathbf{R}}
\providecommand{\EffectResist}[2]{\mathbf{R}_{#1#2}}
\providecommand{\NetworkE}{\EuScript{N}}
\providecommand{\trinom}[4]{\binom{#1}{#2 \;\; #3 \;\; #4}}
\providecommand{\Ranges}{\EuScript{R}}
\providecommand{\CClass}[1]{\ensuremath{\mathbf{#1}}}
\providecommand{\coCClass}[1]{\ensuremath{\mathrm{co-}}\CClass{#1}}
\providecommand{\CX}{\EuScript{C}}
\providecommand{\packet}{\mathbf{c}}
\providecommand{\packetA}{\mathbf{d}}
\providecommand{\mopt}{m_{\mathrm{opt}}}
\providecommand{\girth}{\mathop{\mathrm{girth}}}
\providecommand{\CR}{\mathop{\mathrm{cr}}}
\providecommand{\TRUE}{{\textcolor[named]{RedViolet}{\textsf{TRUE}}}\xspace}
\providecommand{\FALSE}{{\textcolor[named]{RedViolet}{\textsf{FALSE}}}\xspace}
\providecommand{\Part}[1]{\part{#1}}
\providecommand{\Section}[1]{\section{#1}}
\providecommand{\SubSection}[1]{\subsection{#1}}
\providecommand{\SubSubSection}[1]{\subsubsection{#1}}
\providecommand{\Exercise}[5]{%
\begin{exercise_h}[#1]%
   \exelab{#3}{}%
   #4%
\end{exercise_h}%
}
\providecommand{\si}[1]{#1}
\providecommand{\myParagraph}[1]{\noindent\textbf{#1}}
\providecommand{\DefName}[1]{#1}
\providecommand{\deflabel}[1]{#1}
\providecommand{\Description}[1]{}

%-----------------------------------------------------------------------------
% Mathematical Notations & Environments (from notations.sty)
%-----------------------------------------------------------------------------
%-----------------------------------------------------------------------------
% List environments (questions & compactenum)
%-----------------------------------------------------------------------------

\newlist{compactenumA}{enumerate}{5}%
\setlist[compactenumA]{topsep=0pt,itemsep=-1ex,partopsep=1ex,parsep=1ex,%
   label=(\Alph*)}%

\newlist{compactenumI}{enumerate}{5}%
\setlist[compactenumI]{topsep=0pt,itemsep=-1ex,partopsep=1ex,parsep=1ex,%
   label=(\Roman*)}%

\newlist{compactenuma}{enumerate}{5}%
\setlist[compactenuma]{topsep=0pt,itemsep=-1ex,partopsep=1ex,parsep=1ex,%
   label=(\alph*)}%

\newlist{compactitem}{enumerate}{5}%
\setlist[compactitem]{topsep=0pt,itemsep=-1ex,partopsep=1ex,parsep=1ex,%
   label=\textbullet}%

\newlist{compactenumi}{enumerate}{5}%
\setlist[compactenumi]{topsep=0pt,itemsep=-1ex,partopsep=1ex,parsep=1ex,%
   label=(\roman*)}%

\setcounter{secnumdepth}{5}

\newcommand{\Term}[1]{\textsf{#1}}
\newcommand{\TermI}[1]{\Term{#1}\index{#1@\Term{#1}}}

\DeclareMathAlphabet{\mathpzc}{OT1}{pzc}{m}{it}

\newcommand{\seg}{\mathsf{s}}

\newcommand{\Diam}{\mathrm{diam}}
\newcommand{\diam}{\mathrm{diam}}

\newcommand{\Copt}{C_{\mathrm{opt}}}
\newcommand{\CenterPnt}{\mathrm{cen}}

\newcommand{\parent}{\overline{\mathrm{p}}}
\newcommand{\parentX}[2][\!]{\parent\pth[#1]{#2}}

\newcommand{\FastCut}{\Algorithm{FastCut}\xspace}
\newcommand{\MinCut}{\Algorithm{MinCut}\xspace}
\newcommand{\MinCutRep}{\Algorithm{MinCutRep}\xspace}
\newcommand{\Contract}{\Algorithm{Contract}\xspace}

\newcommand{\Grid}{\mathrm{G}}
\newcommand{\ClosestPair}{\mathcal{CP}}
\newcommand{\SphereChar}{\ensuremath{\mathbb{S}}}
\newcommand{\Sphere}[1]{{\SphereChar^{(#1)}}}

\newcommand{\PntSet}{{\mathsf{P}}}%
\providecommand{\PntSetA}{{\mathsf{Q}}}%

\newcommand{\query}{\mathcalb{q}}

\newcommand{\pnt}{\mathsf{p}}
\newcommand{\pntA}{\mathsf{u}}
\newcommand{\pntB}{\mathsf{v}}
\newcommand{\pntC}{{\mathsf{x}}}
\newcommand{\pntD}{{\mathsf{y}}}
\newcommand{\pntE}{{\mathsf{z}}}

\newcommand{\pntCoord}{p}
\newcommand{\pntCoordA}{q}

%%%%%%%%%%%%%%%%%%%%%%%%%%%%%%%%%%%%%%%%%%%%%%%%%%%%%%%%%%%%%%%%%%%%% 
%%%%%%%%%%%%%%%%%%%%%%%%%%%%%%%%%%%%%%%%%%%%%%%%%%%%%%%%%%%%%%%%%%%%%
% Markov chain notations.

\newcommand{\MCSpc}{\mathtt{S}}
\newcommand{\MCPrb}{\mathtt{P}}
\newcommand{\MCPrbY}[2]{\MCPrb_{#1#2}}
\newcommand{\MCPrbM}{\mathbf{P}}%
\newcommand{\MCPrbT}[3]{\MCPrb_{#1#2}^{(#3)}}
\newcommand{\MCFPrbT}[3]{\mathtt{r}_{#1#2}^{(#3)}}
\newcommand{\MCVstProb}[2]{\mathtt{f}_{#1#2}}
\newcommand{\MCExSteps}[2]{\mathtt{h}_{#1#2}}
\newcommand{\MCPVec}[1]{\mathbf{q}^{(#1)}}
%%%%%%%%%%%%%%%%%%%%%%%%%%%%%%%%%%%%%%%%%%%%%%%%%%%%%%%%%%%%%%%%%%%%%

\newcommand{\Spread}{\Phi}
\newcommand{\VC}{\ensuremath{\mathsf{VC}}\xspace}
\newcommand{\MC}{\ensuremath{\mathsf{MC}}\xspace}
\newcommand{\hyperplane}{h}

\newcommand{\QuickSort}{\Algorithm{QuickSort}\xspace}
\newcommand{\Alg}{\Algorithm{Alg}\xspace}

\newcommand{\Graph}{\mathsf{G}}%
\newcommand{\GraphB}{\mathsf{H}}
\newcommand{\GraphC}{\mathsf{K}}
\newcommand{\GraphD}{\mathsf{F}}

\providecommand{\G}{}
\renewcommand{\G}{\Graph}

\newcommand{\RLP}{\ComplexityClass{RLP}\xspace}
\newcommand{\USTCON}{\ComplexityClass{USTCON}}
\newcommand{\STCON}{\overrightarrow{\ComplexityClass{STCON}}}

\newcommand{\Mtrx}{\mathsf{M}}
\newcommand{\TMtrx}{\mathsf{Q}}
\newcommand{\IMtrx}{\mathcal{I}}

\newcommand{\Expander}[3]{$\left[#1,#2,#3\right]$-expander}

\newcommand{\ChartX}[1]{\ensuremath{\chi_{#1}}}
\newcommand{\Field}{\mathbb{F}}
\newcommand{\LSpace}{\mathbb{G}}
\newcommand{\LD}[2]{\mathsf{L{D}}\pth{#1,#2}}
\newcommand{\SEVN}{\widehat{\lambda_2}}
\newcommand{\NX}[2]{#1\!\left[#2\right]}
\newcommand{\Entry}[2]{#1\!\left[#2\right]}
\newcommand{\NXExt}[3]{#1_{#2}\!\left[#3\right]}
\newcommand{\RangeI}[1]{\llbracket #1 \rrbracket}
\newcommand{\DialX}[1]{\gamma\pth{#1}}
\newcommand{\BiDialX}[1]{{\gamma}_2\pth{#1}}

\newcommand{\EV}[1]{\widehat{\lambda_{#1}}}
\newcommand{\TEV}[1]{\widehat{\lambda_{#1}}}
\newcommand{\EVal}{\widehat{\lambda}}
\newcommand{\Basis}{\mathcal{B}}

\newcommand{\ZProd}[2]{\ensuremath{#1 \, \text{\textcircled{z}} #2}}
\newcommand{\RProd}[2]{\ensuremath{#1 \, \text{\textcircled{r}} #2}}

\newcommand{\OneVec}[1]{\mathbf{1}^{#1}}
\newcommand{\DstrbV}[1]{\mathbf{\mathsf{q}}^{(#1)}}
\newcommand{\DstrbE}[2]{\mathsf{q}^{(#1)}_{#2}}

\newcommand{\DEdges}{\widetilde{\Edges}}

\newcommand{\VV}{\mathsf{V}}%
\newcommand{\EE}{\mathsf{E}}%

\newcommand{\Vertices}{\VV}%
\newcommand{\Edges}{\EE}%

\newcommand{\EdgesX}[1]{\Edges\pth{#1}}
\newcommand{\VerticesX}[1]{V\pth{#1}}

\newcommand{\StringX}[1]{\mathsf{s}\pth{#1}}

\newcommand{\BMtrx}{\mathsf{B}}
\newcommand{\WMtrx}{\mathsf{W}}
\newcommand{\NWMtrx}{\overline{\mathsf{W}}}
\newcommand{\Sequence}{\EuScript{S}}

\newcommand{\med}{\mathrm{med}}
\newcommand{\medX}[1]{\mathrm{med}\pth{#1}}

\newcommand{\medianC}{\nu}%

\newlength{\savedparindent}
\newcommand{\SaveIndent}{\setlength{\savedparindent}{\parindent}}
\newcommand{\RestoreIndent}{\setlength{\parindent}{\savedparindent}}

\newcommand{\NNN}{\mathbb{N}}%
\providecommand{\NN}{\mathbb{N}}%
\newcommand{\reals}{\Re}%
\newcommand{\ZZ}{\mathbb{Z}}%

\newcommand{\level}{\mathop{\mathrm{level}}}

\providecommand{\TPDF}[2]{\texorpdfstring{#1}{#2}}

\newcommand{\Hint}[1]{\emph{[#1]}}
\newcommand{\Arr}{\mathop{\mathrm{\EuScript{A}}}}

\newcommand{\Array}{\mathcal{T}}%

\newcommand{\dirEdge}[2]{\left({{#1} \rightarrow {#2}} \right)}

%%%%%%%%%%%%%%%%%%%%%%%%%%%%%%%%%%%%%%%%%%%%%%%%%%%%%%%%%%%%%%%%%%%%%%
% Emphasize stuff...
\definecolor{blue25}{rgb}{0, 0, 11}
\newcommand{\emphic}[2]{%
   \textcolor{blue25}{%
      \bfseries{\emph{#1}}}%
   \index{#2}}
\newcommand{\emphi}[1]{\emphic{#1}{#1}}
\newcommand{\emphOnly}[1]{\textcolor{blue25}{\bfseries{\emph{#1}}}}
\providecommand{\emphiC}[1]{\emphic{#1}{#1}}

\definecolor{almostblack}{rgb}{0, 0, 0.3}
\providecommand{\emphw}[1]{{\textcolor{almostblack}{\emph{#1}}}}%

%
%%%%%%%%%%%%%%%%%%%%%%%%%%%%%%%%%%%%%%%%%%%%%%%%%%%%%%%%%%%%%%%%%%%%%%

\newcommand{\volX}[1]{\mathrm{vol}\pth{#1}}

\newcommand{\Interval}{\EuScript{I}}
\newcommand{\IntervalA}{\EuScript{I}}

\newcommand{\Exp}{\mathrm{Exp}}

\let\Ball\undefined
\newcommand{\Ball}{\mathbf{b}}
\newcommand{\hdim}{d}
\newcommand{\Domain}{\mathcal{D}}
\newcommand{\points}[1]{[\bfseries{#1 points}]}
\newcommand{\Event}{\EuScript{E}}%
\newcommand{\EventB}{\EuScript{B}}%
\newcommand{\EventF}{\EuScript{F}}%
\newcommand{\EventG}{\EuScript{G}}%

\newcommand{\distCmd}[1]{\left\| {#1}  \right\|}
\newcommand{\distY}[2]{\distCmd{#1 - #2}}

\newcommand{\NormalC}{\mathbf{N}}
\newcommand{\NormalY}[2]{\NormalC\pth{#1, #2}}%
\newcommand{\NormalDistOneDim}{\NormalC}
\newcommand{\NormalDistX}[1]{\NormalC^{#1}}
\DeclareMathOperator{\lca}{lca}
\newcommand{\RangeSpace}{\mathsf{S}}%
\newcommand{\RangeSpaceA}{\mathsf{T}}

\newcommand{\range}{\mathbf{r}}
\newcommand{\rangeA}{\mathbf{r}'}

\newcommand{\BadProb}{\varphi}

\newcommand{\EntropyChar}{\mathbb{H}}
\newcommand{\Entropy}[1]{\EntropyChar\pth{#1}}
\newcommand{\EntropyX}[1]{\EntropyChar\pth[]{#1}}

\newcommand{\Matrix}{\mathsf{M}}
\newcommand{\MatrixB}{\mathsf{B}}
\newcommand{\MatrixC}{\mathsf{C}}
\newcommand{\MatrixD}{\mathsf{D}}

\renewcommand{\th}{th\xspace}

\newcommand{\Family}{\mathcal{F}}
\newcommand{\FamilyB}{\mathcal{G}}
\newcommand{\FamilyC}{\mathcal{H}}
\newcommand{\FMSC}{\EuScript{Y}\index{metric space}}

\newcommand{\median}{\overline{\mathsf{m}}}
\newcommand{\QuickSelect}{\Algorithm{Quick{}Select}\xspace}

\providecommand{\row}{\mathbf{r}}%
\newcommand{\rowc}{{r}}%

\newcommand{\vecb}{\mathbf{b}}

\newcommand{\vecA}{\mathbf{v}}
\newcommand{\vecB}{\mathbf{u}}
\newcommand{\vecRand}{\mathbf{r}}

\newcommand{\vecH}{\mathbf{h}}
\newcommand{\vecV}{\mathbf{v}}

\newcommand{\vecX}{\mathbf{x}}
\newcommand{\vecY}{\mathbf{y}}

\newcommand{\DistBinom}[2]{\mathrm{Bin}\pth{#1,#2}}

\newcommand{\DistGeom}[1]{\mathrm{Geom}\pth{#1}}

\newcommand{\Sample}{\mathsf{R}}

\newcommand{\lcm}{\mathop{\mathrm{lcm}}}
\newcommand{\bdiv}{\mathop{\mathrm{div}}}%
\newcommand{\divides}{\mid}%
\newcommand{\poly}{\mathrm{poly}}%

\newcommand{\Group}{\mathcal{G}} 
\newcommand{\GroupA}{\mathcal{H}} 

\newcommand{\algGCD}{\Algorithm{EuclidGCD}\xspace}%

\newcommand{\gOp}{\times}%
\newcommand{\gIdent}{\mathsf{i}}%

\newcommand{\ordX}[2][\!]{\mathop{\mathrm{o{r}d}}\pth[#1]{#2}}

\newcommand{\lSym}[2]{%
   \left(
       \ds {#1}\divides{#2}%
   \right) %
}
\newcommand{\jSym}[2]{\left\llbracket%
       \ds {#1}\divides{#2}%
       \right\rrbracket%
}
\newcommand{\Line}{\ell}

\newcommand{\LevelChar}{\mathbb{L}}
\newcommand{\AtMostKLevelChar}{\LevelChar_{\leq k}}
\newcommand{\AtMostKLevel}[1]{\LevelChar_{\leq k} \pth{ #1}}

\newcommand{\LineSet}{\mathsf{L}}%
\newcommand{\ArrX}[1]{\EuScript{A}\pth{#1}}%

\newcommand{\ProblemC}[1]{{\textsf{\textcolor{red}{#1}\index{problem!#1}}}}
\DefineNamedColor{named}{OliveGreen} {cmyk}{0.64,0,0.95,0.40}
\providecommand{\ComplexityClass}[1]{{{\textcolor[named]{OliveGreen}{%
            \textsc{#1}}}}}
\providecommand{\NPHard}{{\ComplexityClass{NP-Hard}}%
   \index{NP!hard}\xspace}
\providecommand{\NPComplete}{\ComplexityClass{NP-Complete}%
   \index{NP!complete}\xspace}
\providecommand{\APXHard}{\ComplexityClass{APX-Hard}%
   \index{APX!Hard}\xspace}
\providecommand{\NPC}{\ComplexityClass{NPC}%
   \index{NP!complete}\xspace}
\providecommand{\NP}{\ComplexityClass{NP}%
   \index{NP}\xspace}
\providecommand{\POLYT}{\ComplexityClass{P}\xspace}

\newcommand{\envFrame}[1]{\begin{center}\fbox{\begin{minipage}[c]{0.75\textwidth}{#1}\end{minipage}}\end{center}}

\newcommand{\mysqrt}[1]{\ts\!\!\sqrt{#1}}
\newcommand{\SegSet}{\mathsf{S}}
\newcommand{\ArrVD}[1]{\Arr^{|} \pth{#1}}
\newcommand{\RIC}{\textsf{RIC}\index{random incremental
      construction}\xspace}
\newcommand{\ArrB}{\EuScript{B}}

\newcommand{\TreeOpt}{\EuScript{T}_\mathrm{opt}}
\newcommand{\TreeOptA}{\EuScript{T}_\mathrm{opt}'}

\newcommand{\DepthExt}[2]{\mathrm{depth}_{#1}\pth{#2}}
\newcommand{\depthX}[1]{\mathrm{depth}\pth{#1}}

\newcommand{\TreeA}{\EuScript{U}}
\newcommand{\TreeB}{\EuScript{V}}

\newcommand{\Trap}{\sigma}
\newcommand{\TrapA}{\sigma'}
\newcommand{\TrapB}{\sigma''}

\newcommand{\SampleA}{\mathsf{R}'}

\newcommand{\AllRegions}{\mathcal{T}}
\newcommand{\RegionsSet}[1]{\mathcal{F}\pth{#1}}
\newcommand{\RegionsHeavy}[2]{\mathcal{F}_{\!\geq#1}\pth{#2}}
\newcommand{\NumOfRegionsZero}[1]{\ExChar \!f\pth{#1}}
\newcommand{\NumOfRegions}[2]{\ExChar\! f_{\!\geq #1}\pth[]{#2}}
\newcommand{\ObjSet}{\mathsf{S}}

\newcommand{\WeightX}[2][\!]{\omega \pth[#1]{#2}}
\newcommand{\Edge}{\mathsf{e}}

\newcommand{\Face}{\mathsf{f}}

\newcommand{\DefSet}[1]{D\pth{#1}}
\newcommand{\KillSet}[1]{K\pth{#1}}

\newcommand{\CSProbFull}[4]{
   \,\rho_{\!\!\stackrel{\rule[-.005cm]{0cm}{0.1cm}}{#1,
       #4}}\pth{\MakeSBig #2,#3}}

\newcommand{\disk}{D}
\newcommand{\diskA}{D'}

\newcommand{\IntervalSet}{\widehat{\mathcal{I}}}

\newcommand{\GroundSet}{\textsf{X}}
\newcommand{\FGroundSet}{\mathsf{x}}
\newcommand{\RangeSet}{{\mathcal{R}}}
\newcommand{\RangeSetA}{\mathcal{R}'}
\newcommand{\MeasureChar}{\overline{m}}
\newcommand{\Measure}[1]{\MeasureChar\pth{#1}}
\newcommand{\sMeasure}[1]{\overline{s}\pth{#1}}

\newcommand{\SetA}{S}%
\newcommand{\SetB}{T}%
\newcommand{\SetC}{N}%

\newcommand{\VCProj}[2]{#1_{|#2}}
\newcommand{\VCDim}{\mathrm{dim}_{\VC}}
\newcommand{\Dim}{\delta}
\newcommand{\DimA}{\delta'}
\newcommand{\PntSetB}{{{\mathsf{Q}}}}
\newcommand{\PntSetC}{{{\mathsf{R}}}}
\newcommand{\Growth}[2]{{\EuScript{G}}_{#1}\pth{#2}}
\newcommand{\ShatterFunc}[1]{\pi_{#1}}
\newcommand{\Distribution}{\EuScript{D}}
\newcommand{\Error}{\EuScript{E}}

\newcommand{\bd}{{\partial}}%
\newcommand{\eps}{{\varepsilon}}%

\newcommand{\clmlab}[1]{\label{claim:#1}}
\newcommand{\clmref}[1]{\HLink{Claim}{claim:#1}}

\newcommand{\CHX}[1]{{\mathcal{CH}}\pth{#1}}

\newcommand{\ts}{\hspace{0.6pt}}

\def\qedsymbol{\myqedsymbol}%

\newcommand{\qedsign}{\hfill{\hfill\myqedsymbol}}
\newcommand{\qedsignBig}{\hfill{\hfill\rule{4mm}{4mm}}}
\newcommand{\QED}{{\hfill{\hfill\qedsymbol}}}

\newcommand{\num}{\alpha}

\newcommand{\NEW}[1]{#1}

\newcommand{\IncludeGraphics}[2][]{%
   \typeout{}%
   \typeout{Graphics: #2}%
   \typeout{\ includegraphics[#1]{#2}}%
   \includegraphics[#1]{#2}
   \typeout{}%
}

%    -----------------------------------------------------------------
%    -----------------------------------------------------------------
%    -----------------------------------------------------------------
%    Handling references
%    -----------------------------------------------------------------
%    -----------------------------------------------------------------
%    -----------------------------------------------------------------
\newcommand{\HLinkShort}[2]{\hyperref[#2]{#1\ref*{#2}}}
\newcommand{\HLink}[2]{\hyperref[#2]{#1~\ref*{#2}}}
\newcommand{\HLinkPage}[2]{\hyperref[#2]{#1~\ref*{#2}%
      $_\text{p\pageref{#2}}$}}

\newcommand{\HLinkSuffix}[3]{\hyperref[#2]{#1\ref*{#2}{#3}}}
\newcommand{\HLinkPageSuffix}[3]{\hyperref[#2]{#1\ref*{#2}%
      #3$_\text{p\pageref{#2}}$}}

\newcommand{\figlab}[1]{\label{fig:#1}}
\newcommand{\figref}[1]{\HLink{Figure}{fig:#1}}

\newcommand{\exelab}[1]{\label{exer:#1}}

\providecommand{\deflab}[1]{\label{def:#1}}
\newcommand{\defref}[1]{\HLink{Definition}{def:#1}}
\newcommand{\defrefY}[2]{\hyperref[def:#2]{#1}}
\newcommand{\defrefpage}[1]{\HLinkPage{Definition}{def:#1}}

\newcommand{\etal}{\textit{et~al.}\xspace}

\newcommand{\obslab}[1]{\label{observation:#1}}
\newcommand{\obsref}[1]{\HLink{Observation}{observation:#1}}
\newcommand{\obsrefY}[2]{\hyperref[observation:#2]{#1}}

\newcommand{\exmlab}[1]{\label{example:#1}}
\newcommand{\exmref}[1]{\HLink{Example}{example:#1}}

\newcommand{\remlab}[1]{\label{rem:#1}}
\newcommand{\remref}[1]{\HLink{Remark}{rem:#1}}%

\newcommand{\chaplab}[1]{\label{chap:#1}}

\newcommand{\itemlab}[1]{\label{item:#1}}
\newcommand{\itemref}[1]{\HLinkSuffix{}{item:#1}{}}

\newcommand{\tedlab}[1]{\label{tedium:#1}}
\newcommand{\tedref}[1]{\HLink{Tedium}{tedium:#1}}%

% Deferred tedium: `by \TediumHere{...}, we have ...` inside a proof.
% Use \TediumHereY{...}{label} to allow later citations with \tedref{label}.
% Bodies may contain paragraphs and displays. They are typeset, in call order,
% after the enclosing proof, using ordinary tedium numbering and references.
% As with other forward references, compile twice to resolve the links.
% Define macros needed by a body outside the proof: its local scope has ended
% when the deferred body is typeset. Verbatim material is not supported.
\ExplSyntaxOn
\int_new:N \g__notes_tedium_proof_int
\int_new:N \g__notes_tedium_item_int
\seq_new:N \g__notes_tedium_stack_seq
\tl_new:N \l__notes_tedium_proof_tl

\cs_new_protected:Npn \__notes_tedium_begin:
  {
    \int_gincr:N \g__notes_tedium_proof_int
    \seq_new:c
      { g__notes_tedium_\int_use:N \g__notes_tedium_proof_int _seq }
    \seq_gput_right:Nx \g__notes_tedium_stack_seq
      { \int_use:N \g__notes_tedium_proof_int }
  }

\cs_new_protected:Npn \__notes_tedium_print:n #1
  {
    \seq_use:cn { g__notes_tedium_#1_seq } { }
    \cs_undefine:c { g__notes_tedium_#1_seq }
  }

\cs_new_protected:Npn \__notes_tedium_end:
  {
    \seq_gpop_right:NN \g__notes_tedium_stack_seq
      \l__notes_tedium_proof_tl
    \exp_args:NV \__notes_tedium_print:n \l__notes_tedium_proof_tl
  }

\cs_new_protected:Npn \__notes_tedium_queue:nn #1#2
  {
    \seq_if_empty:NTF \g__notes_tedium_stack_seq
      {
        \PackageError{notes}
          {Deferred~tedium~used~outside~a~proof}
          {Use~\string\TediumHere\space or~\string\TediumHereY
           \space inside~a~proof~environment.}
      }
      {
        \seq_get_right:NN \g__notes_tedium_stack_seq
          \l__notes_tedium_proof_tl
        \seq_gput_right:cx
          { g__notes_tedium_\l__notes_tedium_proof_tl _seq }
          {
            \exp_not:N \begin {tedium}
            \exp_not:N \tedlab
              { \exp_not:n {#2} }
            \exp_not:n {#1}
            \exp_not:N \end {tedium}
          }
        \tedref{#2}
      }
  }

\NewDocumentCommand{\TediumHere}{+m}
  {
    \int_gincr:N \g__notes_tedium_item_int
    \exp_args:Nnx \__notes_tedium_queue:nn {#1}
      { auto-td-\int_use:N \g__notes_tedium_item_int }
  }

\NewDocumentCommand{\TediumHereY}{+m m}
  { \__notes_tedium_queue:nn {#1}{#2} }

% Each proof has its own queue; the stack also keeps nested proofs separate.
\clist_map_inline:nn
  { proof,altproof,proofextin,proofnoqed,proofof,proofExt,proofnoproof }
  {
    \AddToHook{env/#1/before}[notes/tedium]{\__notes_tedium_begin:}
    \AddToHook{env/#1/after}[notes/tedium]{\__notes_tedium_end:}
  }
\ExplSyntaxOff

\newcommand{\factlab}[1]{\label{fact:#1}}%

\newcommand{\alglab}[1]{\label{algorithm:#1}}

\newcommand{\net}{\mathcal{N}}%

\newcommand{\lemlab}[1]{\label{lemma:#1}}
\newcommand{\lemref}[1]{\HLink{Lemma}{lemma:#1}}%
\newcommand{\lemrefY}[2]{\hyperref[lemma:#2]{#1}}
\newcommand{\lemrefpage}[1]{\HLinkPage{Lemma}{lemma:#1}}
\newcommand{\lemrefshort}[1]{\HLinkShort{L}{lemma:#1}}

\newcommand{\problab}[1]{\label{prob:#1}}%
\newcommand{\probref}[1]{\HLink{Problem}{prob:#1}}%

\newcommand{\eqlab}[1]{\label{equation:#1}}
\newcommand{\Eqrefbare}[1]{\HLinkSuffix{(}{equation:#1}{)}}
\newcommand{\Eqref}[1]{\HLinkSuffix{Eq.~(}{equation:#1}{)}}
\newcommand{\Eqrefpage}[1]{\HLinkPageSuffix{Eq.~(}{equation:#1}{)}}

\newcommand{\seclab}[1]{\label{sec:#1}}
\newcommand{\secref}[1]{\HLink{Section}{sec:#1}}
\newcommand{\secrefY}[2]{\hyperref[sec:#2]{#1}}
\newcommand{\secrefpage}[1]{\HLinkPage{Section}{sec:#1}}

\newcommand{\thmlab}[1]{{\label{theo:#1}}}
\newcommand{\thmref}[1]{\HLink{Theorem}{theo:#1}}
\newcommand{\thmrefY}[2]{\hyperref[theo:#2]{#1}}
\newcommand{\thmrefpage}[1]{\HLinkPage{Theorem}{theo:#1}}

\newcommand{\corlab}[1]{\label{cor:#1}}
\newcommand{\corref}[1]{\HLink{Corollary}{cor:#1}}%
\newcommand{\correfY}[2]{\hyperref[cor:#2]{#1}}

\newcommand{\Yend}{Y_{\mathrm{end}}}%

%    -----------------------------------------------------------------
%    -----------------------------------------------------------------
%    -----------------------------------------------------------------
%    Handling references
%    -----------------------------------------------------------------
%    -----------------------------------------------------------------
%    -----------------------------------------------------------------

\newcommand{\ceil}[1]{\left\lceil {#1} \right\rceil}
\newcommand{\floor}[1]{\left\lfloor {#1} \right\rfloor}
\newcommand{\brc}[1]{\left\{ {#1} \right\}}
\newcommand{\pth}[2][\!]{\mleft({#2}\mright)}%
\newcommand{\brk}[1]{\mleft[{#1}\mright]}%
\newcommand{\pbrc}[2][\!\!]{#1\left[ {#2} \bigr. \right]}
\newcommand{\pbrcS}[1]{\left[ {#1} \right]}

\newcommand{\dotprod}[2]{{#1}\cdot{#2}}
\newcommand{\DotProd}[2]{\permut{{#1},{#2}}}
\newcommand{\permut}[1]{\left\langle {#1} \right\rangle}

\newcommand{\MakeBig}{}
\newcommand{\MakesBig}{}
\newcommand{\MakeSBig}{}
\newcommand{\MakeVBig}{}

\newcommand{\DistChar}{\mathsf{d}}%
\newcommand{\DistSetY}[2]{\DistChar\pth{#1,#2}}
\newcommand{\distG}[2]{\DistChar_\Graph\pth{#1, #2}}
\newcommand{\DistMtrCharX}[1]{\DistChar_{#1}}

\newcommand{\dY}[2]{\DistChar\pth{#1, #2}}%
\newcommand{\dGY}[2]{\DistChar_\Graph\pth{#1, #2}}%

\newcommand{\cset}{\mathcal{C}}%

\newcommand{\MTR}{\mathcal{M}}
\newcommand{\DistMtr}[2]{\DistChar\pth{#1,#2}}

\newcommand{\dist}[1]{\left\| {#1} \right\|}
\newcommand{\norm}[1]{\left\| {#1} \right\|}

\providecommand{\ds}{\displaystyle}
\newcommand{\sep}[1]{\,\left|\, {#1} \bigr.\right.}%
\newcommand{\sepw}[1]{\left|\, {#1} \right.}%

\newcommand{\ProbExt}[2]{\mathop{\mathbf{Pr}}_{#1}\pbrcx{#2}}
\newcommand{\ProbLTR}{\mathbb{P}}%

\newcommand{\Prob}[1]{\mathop{\ProbLTR} \mleft[ #1 \mright]}%
\newcommand{\ProbCond}[2]{\mathop{\ProbLTR}\!\left[%
       #1 \;\middle\vert\; #2 \right]}
\newcommand{\ProbY}[2]{\ProbLTR^{}_{\!#1}\pbrc{#2}}

\newcommand{\ExChar}{\mathbb{E}}%
\newcommand{\ExSym}{\mathop{\ExChar}}%
\newcommand{\Ex}[2][\!]{\ExSym#1\pbrcx{#2}}
\newcommand{\ExExt}[2]{\ExSym_{#1}\!\pbrcx{#2}}
\newcommand{\ExExtS}[2]{\ExSym_{\substack{#1}}\!\pbrcx{#2}}
\newcommand{\ExCond}[2]{\ExSym\!\left[%
       #1 \;\middle\vert\; #2 \right]}

\newcommand{\VarSymb}{\mathop{\mathbb{V}}}
\newcommand{\Var}[1]{\mathop{\mathbb{V}}\!\pbrcx{#1}}

\newcommand{\pbrcx}[1]{\left[ {#1} \right]}%
\newcommand{\nfrac}[2]{{#1}/{#2}}
\newcommand{\cardin}[1]{\left| {#1} \right|}%

\newcommand\CCancel[2][black]{{\renewcommand\CancelColor{\color{#1}}\cancel{#2}}}

\newcommand{\IncGraphPage}[4][]{%
   \IncludeGraphics[page=#4,#1]{\File{#2/#3}}%
}

\newenvironment{enumEx}{
   \begin{compactenumA}
    }{\end{compactenumA}}

\newcommand{\cl}{{\mathrm{cl}}}

\newcommand{\HSetY}[2]{H\pth{#1, #2}}%
\newcommand{\CHSetY}[2]{C\pth{#1, #2}}%
\newcommand{\OldOm}{{\mathcal{O}}}
\newcommand{\NewOm}{{\mathcal{N}}}

%%%%%%%%%%%%%%%%%%%%%%%%%%%%%%%%%%%%%%%%%%%%%%%%%%%%%%%%%%%%%%%%%%%%%%%
%---------------------------------------------------------------
\newcommand{\IntRange}[1]{\left\llbracket #1 \right\rrbracket}
\newcommand{\IRX}[1]{\left\llbracket #1 \right\rrbracket}%
\newcommand{\IRY}[2]{\left\llbracket #1:#2 \right\rrbracket}

%---------------------------------------------------------------
%%%%%%%%%%%%%%%%%%%%%%%%%%%%%%%%%%%%%%%%%%%%%%%%%%%%%%%%%%%%%%%%%%%%%%%

\newcommand{\obj}{f}%
\newcommand{\Wg}{\omega}%
\newcommand{\TotalW}{W}%

\newcommand{\rankY}[2]{{#1}_{ \permut{#2}}}%%
\newcommand{\areaX}[1]{\mathrm{area}\pth{#1}}%%

\newcommand{\mybinom}[2]%
{\Bigl(\begin{array}{@{}c@{}}#1\\#2\end{array}\Bigr)}

\providecommand{\Set}[2]{\left\{ #1 \;\middle\vert\; #2 \right\}}
\newcommand*{\biga}[1]{\vcenter{\hbox{\scalebox{1.5}{\ensuremath#1}}}}
\newcommand{\CNF}{\TermI{CNF}\xspace}%
\newcommand{\Tree}{T}%

\DeclareSymbolFont{sfletters}{OML}{cmbrm}{m}{it}

\DeclareMathSymbol{\salpha}{\mathord}{sfletters}{"0B}

%%%%%%%%%%%%%%%%%%%%%%%%%%%%%%%%%%%%%%%%%%%%%%%%%%%%%%%%%%%%%%%%%%%%

\providecommand{\ProblemInd}[2]{\textsf{\textcolor{red}{#1}\index{Problem!#2}}}
\newcommand{\NPProblemInd}[4]{
   \begin{center}
      \vspace{0.6cm}
       \fbox{
          \begin{minipage}{0.9\linewidth}
              \vspace{-0.6cm} \hspace{-0.35cm}
              {\large \bfseries \ProblemInd{#1}{#2}}
              \vspace{+0.2cm}%

              \noindent{\textcolor{blue}{\bfseries \textsf{Instance}}:} {#3}
              
              \noindent{\textcolor{blue}{\bfseries \textsf{Question}:}} {#4}
          \end{minipage}
       }
   \end{center}
   \smallskip
   }

\newcommand{\NPProblemOpt}[4]{\NPProblemInd{#1}{#2}{#3}{#4}}

%%%%%%%%%%%%%%%%%%%%%%%%%%%%%%%%%%%%%%%%%%%%%%%%%%%%%%%%%%%%%%%%%%%%

\newcommand{\Hastad}{H{\aa}stad\xspace}

\newcommand{\ANNbr}{\Term{ANN}\index{approximate nearest neighbor}\xspace}
\newcommand{\LSH}{\TermI{L{S}H}\xspace}

\newcommand{\DS}{\EuScript{D}} % DataStructure
\newcommand{\DSC}{\mathsf{D}} % DataStructure

\newcommand{\distH}[2]{\DistChar_{H}\pth{#1, #2}}
\newcommand{\distHY}[2]{\DistChar_{H}\pth{#1, #2}}

\newcommand{\DSAprxNearNbr}{\ensuremath{\DataStructure{{\approx%
         \text{Near{}}}}}%
   \index{near neighbor!data-structure!approximate}%
}
\newcommand{\DataStructure}[1]{\ensuremath{\DS_{#1}}}
\newcommand{\DistNN}[2]{\DistChar\pth{#1,#2}}
\newcommand{\pntc}{{p}}%

\newcommand{\seq}{{M}}%
\newcommand{\cTimes}{\beta}%

\newcommand{\DSTimes}{L}%
\newcommand{\rad}{r}%

\newcommand{\HypercubeX}[1]{\Hypercube^{#1}}
\newcommand{\kFunctions}[2]{\mathrm{combine}\pth{#1,#2}}
\newcommand{\WEqualChar}{\Bumpeq}
\newcommand{\WEqual}[2]{{#1}\WEqualChar{#2}}
\newcommand{\FamilyAmp}{\FamilyC_{\WEqualChar}}
\newcommand{\PSPACE}{\ComplexityClass{PSPACE}\xspace}%

\newcommand{\PS}{P}%

\newcommand{\PSetFar}{\ensuremath{\PntSet^{}_{\geq}\xspace}}%
\newcommand{\DA}{D}%
\newcommand{\uniqX}[1]{\mathrm{uniq}\pth{#1}}%
\newcommand{\diffX}[1]{\mathop{}\!\mathrm{d}#1}

\newcommand{\Frq}{F}%
\newcommand{\frq}{f}%
\newcommand{\STRM}{\mathcal{S}}%
\newcommand{\strm}{{s}}%

\newcommand{\Calinescu}{C\u{a}linescu\xspace}

\newcommand{\fGrowX}[1]{G^{}_{\Dim}\pth{#1}}
\newcommand{\CompressFunc}{\mathtt{Compress}}
\newcommand{\RSA}{\Term{RSA}\xspace}
\newcommand{\MST}{\Term{MST}\xspace}
\providecommand{\CClassPP}{\CClass{P{}P}\xspace}
\providecommand{\ID}{\mathrm{id}}

\newcommand{\EncodeFunc}{\text{\ttfamily{Enc}}}
\newcommand{\DecodeFunc}{\text{\ttfamily{Dec}}}
\newcommand{\discX}[1]{\mathrm{disc}\pth{#1}}%

\newcommand{\propertyC}{\psi}%
\newcommand{\propertyX}[1]{\propertyC\pth{#1}}%
\newcommand{\numC}{\Gamma}%

\newcommand{\degreeX}[1]{\mathrm{degree}\pth{#1}}

\newcommand{\ind}{m}

\providecommand{\nequiv}{\not\equiv}
\providecommand{\sqcupplus}{\sqcup^+}

\newcommand{\AlgorithmI}[1]{{\protect\color{RedViolet}{#1}}}

\newcommand{\binsY}[2]{\sqcup_{\geq #1}\pth{#2}}

\newcommand{\BLX}[1]{{{\sqcupplus{}}_{\geq #1}}}

\newcommand{\pdf}{\textsf{pdf}\xspace}

\newcommand{\cdf}{\textsf{cdf}\xspace}

\newcommand{\costX}[1]{\mathrm{cost}\pth{#1}}

\newcommand{\codeX}[1]{\mathrm{code}\pth{#1}}

\newcommand{\lengthX}[1]{\mathrm{l{}e{}n}\pth{#1}}
